% Source: activation_functions.tex
\begin{notechapter}
\chapter{Desirable Properties of Activation Functions}\label{ch:activation-properties}

In neural networks, activation functions play a crucial 
role in introducing non-linearity, enabling the network 
to learn complex patterns. The properties an activation 
function should ideally have include:

\section{Non-linearity}
The activation function should be non-linear to allow 
the network to model complex relationships between 
inputs and outputs. Without non-linearity, a neural 
network would simply be a linear transformation, no 
matter how many layers it has.

\section{Differentiability}
The activation function should be differentiable (at 
least piecewise) to support gradient-based optimization 
methods like backpropagation. This allows the network 
to compute gradients and update weights during training.

\section{Continuity}
The activation function should be continuous to ensure 
smooth gradient descent optimization. Discontinuities 
can cause instability during training.

\section{Monotonicity}
A monotonic activation function ensures that the 
gradient has a consistent sign, which can help in 
convergence during training. However, this is not 
strictly necessary (e.g., Swish and GELU are non-
monotonic but effective).

\section{Bounded Output (Optional)}
Some activation functions (e.g., sigmoid, tanh) have 
bounded outputs, which can help prevent the network's 
outputs from exploding. However, unbounded functions 
(e.g., ReLU) are also widely used.

\section{Zero-Centered (Optional)}
Activation functions like tanh are zero-centered, 
which can help speed up convergence during training 
by reducing the likelihood of zig-zagging updates in 
gradient descent. However, this is not a strict 
requirement (e.g., ReLU is not zero-centered but 
still effective).

\section{Computational Efficiency}
The activation function should be computationally 
efficient to evaluate, as it is applied to every 
neuron in the network during both forward and 
backward passes.

\section{Avoidance of Vanishing Gradients}
The activation function should not saturate too 
quickly (e.g., sigmoid or tanh for large inputs), 
as this can lead to vanishing gradients, making 
training slow or ineffective. ReLU and its variants 
address this issue.

\section{Avoidance of Exploding Gradients}
The activation function should not produce outputs 
that are too large, as this can lead to exploding 
gradients, causing instability during training.

\section{Common Activation Functions and Their Properties}
\begin{itemize}
    \item \textbf{ReLU (Rectified Linear Unit)}: 
          \( f(x) = \max(0, x) \)
          \begin{itemize}
              \item Non-linear, computationally efficient, 
                    avoids vanishing gradients (for positive 
                    inputs), but not zero-centered and can 
                    cause "dying ReLU" issues.
          \end{itemize}
    \item \textbf{Sigmoid}: 
          \( f(x) = \frac{1}{1 + e^{-x}} \)
          \begin{itemize}
              \item Smooth, bounded, but suffers from 
                    vanishing gradients and is not zero-
                    centered.
          \end{itemize}
    \item \textbf{Tanh}: 
          \( f(x) = \tanh(x) \)
          \begin{itemize}
              \item Zero-centered, bounded, but can suffer 
                    from vanishing gradients.
          \end{itemize}
    \item \textbf{Leaky ReLU}: 
          \( f(x) = \max(0.01x, x) \)
          \begin{itemize}
              \item Addresses the "dying ReLU" problem by 
                    allowing small negative outputs.
          \end{itemize}
    \item \textbf{Softmax}: Used in the output layer for 
          classification tasks to produce probability 
          distributions.
\end{itemize}

Choosing the right activation function depends on the 
specific problem, network architecture, and training 
dynamics.

\chapterappendix
\section{Appendix: The Problem of Non-Zero-Centered Activation Functions}

Non-zero-centered activation functions (e.g., sigmoid, ReLU) can lead to 
inefficient training and zig-zagging updates in gradient descent. This 
appendix explains why this happens and how zero-centered activations 
can help.

\subsection{Non-Zero-Centered Activations}
A non-zero-centered activation function produces outputs that are 
biased in one direction (e.g., always positive for ReLU or always in 
the range \( (0, 1) \) for sigmoid). For example:
\begin{itemize}
    \item \textbf{ReLU}: \( f(x) = \max(0, x) \) (outputs are \( \geq 0 \)),
    \item \textbf{Sigmoid}: \( f(x) = \frac{1}{1 + e^{-x}} \) (outputs are \( \in (0, 1) \)).
\end{itemize}

\subsection{Impact on Gradients}
During backpropagation, the gradient of the loss \( L \) with respect 
to a weight \( w_i \) is:
\[
\frac{\partial L}{\partial w_i} = \frac{\partial L}{\partial a_l} \cdot \frac{\partial a_l}{\partial z} \cdot \frac{\partial z}{\partial w_i},
\]
where:
\begin{itemize}
    \item \( \frac{\partial z}{\partial w_i} = a_{l-1,i} \) (the activation of the previous layer),
    \item \( a_{l-1,i} \) is the input to the current neuron.
\end{itemize}

If the activations \( a_{l-1,i} \) are all positive (due to a non-zero-centered 
activation function), the gradients \( \frac{\partial L}{\partial w_i} \) will 
tend to have the same sign. This causes all weight updates to move in the 
same direction, leading to inefficient, zig-zagging updates.

\subsection{Why Zero-Centered Activations Help}
Zero-centered activation functions (e.g., tanh) produce outputs that are 
symmetrically distributed around zero (both positive and negative). This 
ensures that the activations \( a_{l-1,i} \) can be both positive and negative, 
allowing the gradients \( \frac{\partial L}{\partial w_i} \) to have both 
positive and negative values. As a result, weight updates can move in 
different directions, leading to smoother and more efficient optimization.

\subsection{Practical Implications}
While zero-centered activations like tanh can improve convergence, non-zero-centered 
activations like ReLU are still widely used due to their simplicity and effectiveness. 
Techniques like batch normalization can mitigate the issues caused by non-zero-centered 
activations by normalizing the inputs to each layer.
\end{notechapter}
